\begin{center}
\textsc{Gauss} an \textsc{Bessel}.\quad Göttingen, 12.~März 1826.
\end{center}

\par\noindent\hspace*{\parindent}.\ .\ .\ .\ .\ .\quad Den grössten Theil des vorigen Sommers habe ich im Bremischen und Oldenburgschen mit meinen Messungen zugebracht.\quad.\ .\ .\ .\ .\ .

\par\noindent\hspace*{\parindent}Was meine Messungen betrifft, so habe ich deren trigonometrischen Theil, wenigstens dem Buchstaben nach, vollendet; meine Seite Varel-Jever schliesst sich an die \textsc{Krayenhoff}schen Dreiecke.
Ob sie aber wirklich geendigt sind, weiss ich selbst noch nicht.
Meine Winkelmessungen in Jever geben \textls[120]{ganz enorme} Unterschiede von den \textsc{Krayenhoff}schen, die bis auf \(15^{\prime\prime}\) gehen.
Ihre Quelle
\par\clearpage
\noindent{\small\textbf{ZUR HANNOVERSCHEN TRIANGULATION.}}\hfill361\par\medskip
\noindent kann ich nicht mit Gewissheit angeben.
Meine eigenen Winkel verbürge ich bis auf eine Secunde.
Die Differenzen würden sich erklären lassen, wenn ich annehmen dürfte, dass \textsc{Krayenhoffs} Centrum der Station von dem meinigen (ich habe eine Etage tiefer beobachtet) ein Meter entfernt liegt.
Die Aufschlüsse, die ich von \textsc{Krayenhoff} erhalten habe, sind unbefriedigend in Rücksicht auf das Centriren; im allgemeinen aber geht daraus hervor, dass seine Winkelmessungen in dortiger Gegend lange nicht die Schärfe haben wie die meinigen.
Er hat in Ostfriesland ein schlechteres Instrument gebraucht; aus vielen Winkelreihen hat er immer nur diejenigen beibehalten, die am besten zu passen schienen (ohne anzugeben, wie viel die andern abwichen) und selbst unter den beibehaltenen in Jever finden sich Differenzen von \(4^{\prime\prime}\).
Unter diesen Umständen scheint es mir nicht rathsam, meine Messungen in Ostfriesland weiter auszudehnen.
Seine südlichern Messungen sind besser, und eine Verbindung meiner Dreiecke über das Osnabrücksche nach Bentheim würde ohne Zweifel zuverlässigere Resultate geben können.
Allein der wankende Zustand meiner Gesundheit hat mich bisher muthlos gemacht, auf solche Operationen, die noch eine ein- oder anderthalbjährige Campagne erfordern würden, anzutragen, und in diesem Sommer wird schwerlich etwas erhebliches darin geschehen können.\quad.\ .\ .\ .\ .\ .

